\documentclass[11pt,a4paper]{article}

% ===== PACKAGES =====
\usepackage[utf8]{inputenc}
\usepackage[T1]{fontenc}
\usepackage[margin=1in]{geometry}
\usepackage{hyperref}
\usepackage{listings}
\usepackage{xcolor}
\usepackage{booktabs}
\usepackage{array}
\usepackage{longtable}
\usepackage{fancyhdr}
\usepackage{titlesec}
\usepackage{tcolorbox}
\usepackage{fontawesome5}
\usepackage{amssymb}
\usepackage{amsmath}
\usepackage{enumitem}
\usepackage{graphicx}

% ===== COLORS =====
\definecolor{codegreen}{rgb}{0.25,0.49,0.48}
\definecolor{codegray}{rgb}{0.5,0.5,0.5}
\definecolor{codepurple}{rgb}{0.58,0,0.82}
\definecolor{codeblue}{rgb}{0.13,0.29,0.53}
\definecolor{codeorange}{rgb}{0.8,0.4,0.0}
\definecolor{backcolour}{rgb}{0.95,0.95,0.95}
\definecolor{warningbg}{rgb}{1.0,0.95,0.85}
\definecolor{warningborder}{rgb}{0.9,0.6,0.2}
\definecolor{tipbg}{rgb}{0.9,0.97,0.9}
\definecolor{tipborder}{rgb}{0.3,0.7,0.3}
\definecolor{notebg}{rgb}{0.9,0.93,1.0}
\definecolor{noteborder}{rgb}{0.3,0.4,0.7}

% ===== IEC 61131-3 CODE LISTING STYLE =====
\lstdefinelanguage{IEC61131}{
    keywords={VAR, END_VAR, VAR_INPUT, VAR_OUTPUT, VAR_IN_OUT, VAR_GLOBAL, 
              TYPE, END_TYPE, STRUCT, END_STRUCT, FUNCTION, END_FUNCTION,
              FUNCTION_BLOCK, END_FUNCTION_BLOCK, PROGRAM, END_PROGRAM,
              IF, THEN, ELSE, ELSIF, END_IF, CASE, OF, END_CASE,
              FOR, TO, BY, DO, END_FOR, WHILE, END_WHILE, REPEAT, UNTIL, END_REPEAT,
              TRUE, FALSE, AND, OR, NOT, XOR, MOD, RETURN,
              ARRAY, OF, POINTER, REFERENCE, TO, AT, EXTENDS, IMPLEMENTS,
              METHOD, END_METHOD, PROPERTY, END_PROPERTY,
              INTERFACE, END_INTERFACE,
              PUBLIC, PRIVATE, PROTECTED, INTERNAL, FINAL, ABSTRACT,
              THIS, SUPER, VAR_STAT, CONSTANT, VAR_INST},
    keywordstyle=\color{codeblue}\bfseries,
    keywords=[2]{BOOL, BYTE, WORD, DWORD, LWORD, SINT, USINT, INT, UINT, 
                 DINT, UDINT, LINT, ULINT, REAL, LREAL, STRING, WSTRING,
                 TIME, LTIME, DATE, TIME_OF_DAY, TOD, DATE_AND_TIME, DT,
                 FB_MBReadCoils, FB_MBWriteCoils, FB_MBReadRegs, 
                 FB_MBWriteSingleRegister, FB_MBWriteRegs,
                 E_ReadState, E_WriteState},
    keywordstyle=[2]\color{codepurple}\bfseries,
    keywords=[3]{ADR, SIZEOF, REF, INT_TO_WORD, WORD_TO_INT, BOOL_TO_BYTE},
    keywordstyle=[3]\color{codeorange}\bfseries,
    sensitive=false,
    morecomment=[l]{//},
    morecomment=[s]{(*}{*)},
    commentstyle=\color{codegreen}\itshape,
    stringstyle=\color{codeorange},
    morestring=[b]',
    morestring=[b]",
}

\lstdefinestyle{iecstyle}{
    language=IEC61131,
    backgroundcolor=\color{backcolour},
    basicstyle=\ttfamily\small,
    breakatwhitespace=false,
    breaklines=true,
    captionpos=b,
    keepspaces=true,
    numbers=left,
    numbersep=8pt,
    numberstyle=\tiny\color{codegray},
    showspaces=false,
    showstringspaces=false,
    showtabs=false,
    tabsize=4,
    frame=single,
    framerule=0.5pt,
    rulecolor=\color{codegray},
    xleftmargin=15pt,
    framexleftmargin=15pt,
}

\lstset{style=iecstyle}

% ===== TCOLORBOX STYLES =====
\tcbuselibrary{skins,breakable}

\newtcolorbox{warningbox}{
    colback=warningbg,
    colframe=warningborder,
    boxrule=1pt,
    left=10pt,
    right=10pt,
    top=8pt,
    bottom=8pt,
    breakable,
    before upper={\faExclamationTriangle\ \textbf{Warning:} }
}

\newtcolorbox{tipbox}{
    colback=tipbg,
    colframe=tipborder,
    boxrule=1pt,
    left=10pt,
    right=10pt,
    top=8pt,
    bottom=8pt,
    breakable,
    before upper={\faLightbulb\ \textbf{Tip:} }
}

\newtcolorbox{notebox}{
    colback=notebg,
    colframe=noteborder,
    boxrule=1pt,
    left=10pt,
    right=10pt,
    top=8pt,
    bottom=8pt,
    breakable,
    before upper={\faInfoCircle\ \textbf{Note:} }
}

% ===== HEADER/FOOTER =====
\pagestyle{fancy}
\fancyhf{}
\fancyhead[L]{\leftmark}
\fancyhead[R]{TwinCAT 3 Lecture Notes}
\fancyfoot[C]{\thepage}
\renewcommand{\headrulewidth}{0.4pt}
\renewcommand{\footrulewidth}{0.4pt}

% ===== HYPERREF SETUP =====
\hypersetup{
    colorlinks=true,
    linkcolor=codeblue,
    filecolor=magenta,
    urlcolor=cyan,
    pdftitle={TwinCAT 3 - Extending Functionality with TwinCAT Functions},
    pdfauthor={},
    bookmarks=true,
}

% ===== TITLE FORMATTING =====
\titleformat{\section}
    {\normalfont\Large\bfseries}{\thesection}{1em}{}[{\titlerule[0.8pt]}]
\titleformat{\subsection}
    {\normalfont\large\bfseries}{\thesubsection}{1em}{}
\titleformat{\subsubsection}
    {\normalfont\normalsize\bfseries}{\thesubsubsection}{1em}{}

% ===== DOCUMENT INFO =====
\title{
    \vspace{-1cm}
    \textbf{Lecture Notes: TwinCAT 3}\\
    \Large Extending Functionality with TwinCAT Functions
}
\author{}
\date{\today}

% ===== BEGIN DOCUMENT =====
\begin{document}

\maketitle

\tableofcontents
\newpage

% ============================================================
\section{Introduction to TwinCAT Functions}
% ============================================================

While the base installation of TwinCAT 3 (XAE) covers the majority of industrial automation needs, certain specialized tasks require extending the system's core capabilities. This is achieved through \textbf{TwinCAT Functions}, an ecosystem of extensions categorized by their specific application area.

\begin{notebox}
TwinCAT Functions are add-on modules that extend the base TwinCAT functionality. They are installed separately and require their own licenses (though trial licenses are available for testing).
\end{notebox}

% ============================================================
\section{Categories of TwinCAT Functions}
% ============================================================

TwinCAT functions are organized into several major series:

\begin{table}[h]
\centering
\begin{tabular}{@{}lll@{}}
\toprule
\textbf{Series} & \textbf{Category} & \textbf{Examples} \\
\midrule
TF1xxx & System & Matlab/Simulink, UML editors \\
TF2xxx & HMI & HTML5, CSS, JavaScript interfaces \\
TF3xxx & Measurement & TwinCAT Scope, power measurement \\
TF4xxx & Controller & PID controllers, filters, speech commands \\
TF5xxx & Motion Control & Axis control, multi-axis synchronization \\
TF6xxx & Connectivity & JSON, OPC-UA, Modbus TCP, PROFINET, FTP \\
TF7xxx & Vision & Camera integration, quality inspection \\
TF8xxx & Industry Specific & Building automation, wind power \\
\bottomrule
\end{tabular}
\caption{TwinCAT Function Categories}
\end{table}

\subsection{TF1xxx (System)}
Extensions for the PLC runtime, such as running Matlab/Simulink models or using UML editors.

\subsection{TF2xxx (HMI)}
Tools for creating modern operator interfaces using HTML5, CSS, and JavaScript.

\subsection{TF3xxx (Measurement)}
Tools for analytics, power measurement, and the TwinCAT Scope for process analysis.

\subsection{TF4xxx (Controller)}
Includes PID controllers, various filters, and even speech-to-PLC command functions.

\subsection{TF5xxx (Motion Control)}
A massive category covering all aspects of motion, from simple axis control to complex multi-axis synchronization.

\subsection{TF6xxx (Connectivity)}
Enables the PLC to communicate with external systems via protocols like:
\begin{itemize}[leftmargin=*]
    \item JSON (TF6020)
    \item OPC-UA (TF6100)
    \item Modbus TCP (TF6250)
    \item PROFINET (TF6270)
    \item FTP (TF6300)
\end{itemize}

\subsection{TF7xxx (Vision)}
Integrates camera support directly into the TwinCAT environment for quality inspection.

\subsection{TF8xxx (Industry Specific)}
Specialized frameworks for industries like building automation or wind power.

% ============================================================
\section{Architecture of a Connectivity Function (Modbus TCP)}
% ============================================================

Using a function like \textbf{TF6250 (Modbus TCP)} involves a specific architecture to bridge the gap between real-time control and standard networking.

\subsection{Installation}
The function is a separate executable downloaded from the Beckhoff website and installed on the PLC runtime (XAR).

\subsection{The Process}
Installation creates a background process in the Windows operating system (user-land).

\subsection{Communication Bridge}
The real-time TwinCAT kernel communicates with this Windows process (likely via ADS). The Windows process then utilizes the standard network stack to talk to external Modbus devices.

\subsection{Asynchronous Execution}
Because the networking happens in non-deterministic Windows space, the communication is \textbf{asynchronous} to the real-time PLC cycle.

\begin{figure}[h]
\centering
\begin{tabular}{|c|c|c|c|}
\hline
\textbf{PLC Code} & $\xleftrightarrow{\text{ADS}}$ & \textbf{Windows Process} & $\xleftrightarrow{\text{TCP/IP}}$ \\
(Real-time) & & (TF6250 Service) & \\
\hline
\multicolumn{2}{|c|}{TwinCAT Kernel} & \multicolumn{2}{c|}{Windows User Space} \\
\hline
\end{tabular}

\vspace{0.3cm}

\begin{tabular}{|c|}
\hline
\textbf{Modbus Device} \\
(External PLC, Sensor, etc.) \\
\hline
\end{tabular}
\caption{TF6250 Modbus TCP Architecture}
\end{figure}

% ============================================================
\section{Programming and Licensing}
% ============================================================

\subsection{Libraries}
To use the function, you must add a reference to the specific library (e.g., \texttt{Tc2\_ModbusTCP}) in your PLC project to access its function blocks.

\subsection{Licensing}
These extensions are not free. They require a specific license (e.g., TF6250) added in the ``Manage Licenses'' tab. For testing, 7-day trial licenses can be generated.

\subsection{Modbus Protocol Basics}
Modbus is a primitive protocol from the 1970s that works with bits (coils) and 16-bit values (registers) rather than named variables.

\begin{table}[h]
\centering
\begin{tabular}{@{}llll@{}}
\toprule
\textbf{Modbus Type} & \textbf{Data Type} & \textbf{Access} & \textbf{Function Codes} \\
\midrule
Coils & Single bit (BOOL) & Read/Write & 01, 05, 15 \\
Discrete Inputs & Single bit (BOOL) & Read only & 02 \\
Holding Registers & 16-bit (WORD) & Read/Write & 03, 06, 16 \\
Input Registers & 16-bit (WORD) & Read only & 04 \\
\bottomrule
\end{tabular}
\caption{Modbus Data Types}
\end{table}

\begin{warningbox}
ADS Error \textbf{1828 (0x724)} specifically indicates a missing license for the TwinCAT function. Ensure you have activated the appropriate license or trial before testing.
\end{warningbox}

% ============================================================
\section{Implementation: State Machine Logic}
% ============================================================

Because Modbus communication is asynchronous, you cannot simply call a function block and expect an immediate result. You must implement a state machine to manage the request.

\subsection{State Machine Enumeration}

\begin{lstlisting}
// Local enumeration for managing asynchronous Modbus read operations
// Using local enums prevents namespace pollution and makes code self-documenting
// The {attribute 'strict'} ensures type safety

FUNCTION_BLOCK FB_ModbusReader
VAR
    // Local enumeration for state machine states
    // Defines the asynchronous communication flow
    {attribute 'strict'}
    E_ReadState : (
        WAIT,       // Idle state - waiting for trigger command
        TRIGGER,    // Send the Modbus request (set bExecute := TRUE)
        BUSY        // Wait for response (monitor bBusy and bError)
    );
    
    // Current state variable
    eReadState : E_ReadState := E_ReadState.WAIT;
    
    // Trigger to start a read operation
    bStartRead : BOOL;
    
    // Status outputs
    bReadComplete : BOOL;
    bReadError    : BOOL;
    nErrorCode    : UDINT;
END_VAR
\end{lstlisting}

\begin{lstlisting}
// Alternative: Separate enumeration for write operations
// You may need different states for different operation types

VAR
    // Enumeration for write state machine
    {attribute 'strict'}
    E_WriteState : (
        IDLE,           // Waiting for write command
        SEND_REQUEST,   // Trigger the write operation
        WAIT_RESPONSE,  // Wait for confirmation
        COMPLETE,       // Write successful
        ERROR           // Write failed
    );
    
    eWriteState : E_WriteState := E_WriteState.IDLE;
END_VAR
\end{lstlisting}

\begin{tipbox}
Using local enumerations within function blocks (rather than global types) keeps your code organized and prevents naming conflicts when multiple function blocks need similar state definitions.
\end{tipbox}

\subsection{Modbus Read Coils Implementation}

\begin{lstlisting}
// Modbus TCP Read Coils Implementation
// Reads digital inputs (bits) from a remote Modbus device
// Library required: Tc2_ModbusTCP

FUNCTION_BLOCK FB_ModbusCoilReader
VAR
    // State machine enumeration
    {attribute 'strict'}
    E_ReadState : (
        WAIT,       // Idle - waiting for trigger
        TRIGGER,    // Start the read operation
        READ        // Wait for completion (BUSY state)
    );
    
    eReadState : E_ReadState := E_ReadState.WAIT;
    
    // Modbus Read Coils function block instance
    // From Tc2_ModbusTCP library
    fbReadCoils : FB_MBReadCoils;
    
    // Local buffer to store the read coil data
    // Each BYTE can hold 8 coil values (bits)
    // Size depends on nQuantity: need nQuantity/8 bytes (rounded up)
    aCoilBuffer : ARRAY[0..9] OF BYTE;  // Buffer for up to 80 coils
    
    // Configuration parameters
    sDeviceIP   : STRING(15) := '192.168.1.100';  // Target device IP
    nTCPPort    : UINT := 502;                     // Standard Modbus port
    nUnitID     : BYTE := 255;                     // Modbus unit ID (255 = broadcast)
    nStartAddr  : UINT := 0;                       // Starting coil address
    nNumCoils   : UINT := 8;                       // Number of coils to read
    
    // Control and status
    bExecuteRead   : BOOL;           // Trigger to start read
    bReadDone      : BOOL;           // Read completed successfully
    bReadError     : BOOL;           // Read failed
    nErrorID       : UDINT;          // Error code if failed
    
    // Extracted coil values (for convenience)
    bCoil0         : BOOL;           // First coil value
    bCoil1         : BOOL;           // Second coil value
END_VAR
\end{lstlisting}

\begin{lstlisting}
// State machine implementation for reading coils
CASE eReadState OF

    E_ReadState.WAIT:
        // Idle state - waiting for trigger
        bReadDone := FALSE;
        bReadError := FALSE;
        
        IF bExecuteRead THEN
            bExecuteRead := FALSE;  // Reset trigger
            eReadState := E_ReadState.TRIGGER;
        END_IF
        
    E_ReadState.TRIGGER:
        // Start the Modbus read operation
        // Configure and trigger the function block
        fbReadCoils(
            sIPAddr     := sDeviceIP,           // Target device IP address
            nTCPPort    := nTCPPort,            // TCP port (502 standard)
            nUnitID     := nUnitID,             // Modbus unit identifier
            nQuantity   := nNumCoils,           // Number of coils to read
            nMBAddr     := nStartAddr,          // Starting Modbus address
            cbLength    := SIZEOF(aCoilBuffer), // Buffer size
            pDestAddr   := ADR(aCoilBuffer),    // Pointer to destination buffer
            bExecute    := TRUE,                // TRIGGER the operation
            tTimeout    := T#5S                 // 5 second timeout
        );
        
        eReadState := E_ReadState.READ;
        
    E_ReadState.READ:
        // Wait for operation to complete
        // Must continue calling FB with bExecute := FALSE
        fbReadCoils(bExecute := FALSE);
        
        // Check if operation is complete
        IF NOT fbReadCoils.bBusy THEN
            IF fbReadCoils.bError THEN
                // Read failed
                bReadError := TRUE;
                nErrorID := fbReadCoils.nErrId;
            ELSE
                // Read successful - extract coil values
                bReadDone := TRUE;
                
                // Extract individual bits from buffer
                // Coil 0 is bit 0 of byte 0, etc.
                bCoil0 := (aCoilBuffer[0] AND 16#01) > 0;
                bCoil1 := (aCoilBuffer[0] AND 16#02) > 0;
            END_IF
            
            eReadState := E_ReadState.WAIT;
        END_IF

END_CASE
\end{lstlisting}

\begin{notebox}
The \texttt{ADR()} function returns a pointer to the variable. This is required because the Modbus function blocks write data directly to memory at the specified address. The \texttt{cbLength} parameter tells the function block how much space is available in the buffer.
\end{notebox}

% ============================================================
\section{Practical Workflow: Reading and Writing}
% ============================================================

\begin{enumerate}[leftmargin=*]
    \item \textbf{Wait State:} The program remains idle until a trigger is received.
    \item \textbf{Trigger State:} Set the \texttt{execute} flag of the Modbus function block to \textbf{TRUE} to send the command to the Windows process.
    \item \textbf{Busy/Read State:} Set \texttt{execute} back to \textbf{FALSE} and monitor the \texttt{bBusy} and \texttt{bError} outputs.
    \item \textbf{Completion:} Once \texttt{bBusy} is false, check \texttt{bError}. If successful, the data is available in your local buffer.
\end{enumerate}

\begin{figure}[h]
\centering
\begin{tabular}{|c|c|c|c|}
\hline
\textbf{WAIT} & $\xrightarrow{\text{Trigger}}$ & \textbf{TRIGGER} & $\xrightarrow{\text{Immediate}}$ \\
\hline
\multicolumn{2}{|c|}{bExecute = FALSE} & \multicolumn{2}{c|}{bExecute = TRUE} \\
\hline
\end{tabular}

\vspace{0.3cm}

\begin{tabular}{|c|c|c|}
\hline
$\xrightarrow{\text{Immediate}}$ & \textbf{BUSY/READ} & $\xrightarrow{\text{!bBusy}}$ \\
\hline
\multicolumn{2}{|c|}{bExecute = FALSE, monitor bBusy} & Complete or Error \\
\hline
\end{tabular}
\caption{Asynchronous Communication State Flow}
\end{figure}

\subsection{Modbus Write Register Implementation}

\begin{lstlisting}
// Modbus TCP Write Single Register Implementation
// Writes a 16-bit value to a remote Modbus device
// Library required: Tc2_ModbusTCP

FUNCTION_BLOCK FB_ModbusRegisterWriter
VAR
    // State machine enumeration
    {attribute 'strict'}
    E_WriteState : (
        IDLE,           // Waiting for write command
        SEND,           // Trigger the write operation
        WAIT_COMPLETE   // Wait for confirmation
    );
    
    eWriteState : E_WriteState := E_WriteState.IDLE;
    
    // Modbus Write Single Register function block instance
    // From Tc2_ModbusTCP library
    fbWriteRegister : FB_MBWriteSingleRegister;
    
    // Configuration parameters
    sDeviceIP    : STRING(15) := '192.168.1.100';  // Target device IP
    nTCPPort     : UINT := 502;                     // Standard Modbus port
    nUnitID      : BYTE := 1;                       // Modbus unit ID
    nRegisterAddr: UINT := 100;                     // Register address to write
    
    // Value to write (as INT for user convenience)
    nWriteValue  : INT := 0;
    
    // Internal: Value converted to WORD for Modbus
    wRegisterValue : WORD;
    
    // Control and status
    bExecuteWrite  : BOOL;           // Trigger to start write
    bWriteDone     : BOOL;           // Write completed successfully
    bWriteError    : BOOL;           // Write failed
    nErrorID       : UDINT;          // Error code if failed
END_VAR
\end{lstlisting}

\begin{lstlisting}
// State machine implementation for writing a single register
CASE eWriteState OF

    E_WriteState.IDLE:
        // Idle state - waiting for write command
        bWriteDone := FALSE;
        bWriteError := FALSE;
        
        IF bExecuteWrite THEN
            bExecuteWrite := FALSE;  // Reset trigger
            
            // IMPORTANT: Convert INT to WORD for Modbus protocol
            // Modbus registers are 16-bit unsigned values
            // Use appropriate conversion based on your data type
            wRegisterValue := INT_TO_WORD(nWriteValue);
            
            eWriteState := E_WriteState.SEND;
        END_IF
        
    E_WriteState.SEND:
        // Start the Modbus write operation
        fbWriteRegister(
            sIPAddr     := sDeviceIP,           // Target device IP address
            nTCPPort    := nTCPPort,            // TCP port (502 standard)
            nUnitID     := nUnitID,             // Modbus unit identifier
            nMBAddr     := nRegisterAddr,       // Register address to write
            nValue      := wRegisterValue,      // Value to write (WORD type!)
            bExecute    := TRUE,                // TRIGGER the operation
            tTimeout    := T#5S                 // 5 second timeout
        );
        
        eWriteState := E_WriteState.WAIT_COMPLETE;
        
    E_WriteState.WAIT_COMPLETE:
        // Wait for write operation to complete
        fbWriteRegister(bExecute := FALSE);
        
        IF NOT fbWriteRegister.bBusy THEN
            IF fbWriteRegister.bError THEN
                // Write failed
                bWriteError := TRUE;
                nErrorID := fbWriteRegister.nErrId;
            ELSE
                // Write successful
                bWriteDone := TRUE;
            END_IF
            
            eWriteState := E_WriteState.IDLE;
        END_IF

END_CASE
\end{lstlisting}

\begin{lstlisting}
// Example: Writing multiple registers
// Use FB_MBWriteRegs for writing multiple consecutive registers

VAR
    fbWriteMultiple : FB_MBWriteRegs;
    
    // Buffer containing values to write
    // Each element is a WORD (16-bit)
    aWriteBuffer : ARRAY[0..4] OF WORD := [100, 200, 300, 400, 500];
END_VAR

// In SEND state:
fbWriteMultiple(
    sIPAddr     := '192.168.1.100',
    nTCPPort    := 502,
    nUnitID     := 1,
    nMBAddr     := 100,                    // Starting address
    nQuantity   := 5,                      // Number of registers
    cbLength    := SIZEOF(aWriteBuffer),   // Buffer size in bytes
    pSrcAddr    := ADR(aWriteBuffer),      // Pointer to source data
    bExecute    := TRUE,
    tTimeout    := T#5S
);
\end{lstlisting}

\begin{warningbox}
Modbus uses 16-bit unsigned WORD values for registers. When writing signed integers (INT), you must convert them using \texttt{INT\_TO\_WORD()}. When reading, convert back using \texttt{WORD\_TO\_INT()}. Failure to convert properly can result in incorrect values, especially for negative numbers.
\end{warningbox}

% ============================================================
\section{Key Considerations}
% ============================================================

\subsection{Virtual Machines}
When testing communication between separate systems (e.g., a PLC and a simulator), ensure they have distinct IP addresses on the same network.

\subsection{Error Handling}
Always monitor error codes. Common error codes include:

\begin{table}[h]
\centering
\begin{tabular}{@{}lll@{}}
\toprule
\textbf{Error Code} & \textbf{Hex} & \textbf{Description} \\
\midrule
1828 & 0x724 & Missing TwinCAT function license \\
1861 & 0x745 & Target device not responding \\
6 & 0x06 & Target busy (device specific) \\
274 & 0x112 & Invalid Modbus address \\
\bottomrule
\end{tabular}
\caption{Common Modbus TCP Error Codes}
\end{table}

\subsection{Port 502}
This is the standard TCP port for Modbus communication. Ensure firewalls allow traffic on this port.

\begin{lstlisting}
// Comprehensive error handling example
IF fbReadCoils.bError THEN
    CASE fbReadCoils.nErrId OF
        16#724:  // 1828 decimal
            sErrorMsg := 'License missing for TF6250';
            
        16#745:  // 1861 decimal
            sErrorMsg := 'Device not responding - check IP/connection';
            
        16#112:  // 274 decimal
            sErrorMsg := 'Invalid Modbus address';
            
    ELSE
        sErrorMsg := CONCAT('Unknown error: ', UDINT_TO_STRING(fbReadCoils.nErrId));
    END_CASE
END_IF
\end{lstlisting}

\begin{tipbox}
When testing Modbus communication, use a Modbus simulator (such as ModRSsim2 or Modbus Poll) on a separate machine or VM. This allows you to verify your PLC code is working correctly before connecting to actual field devices.
\end{tipbox}

% ============================================================
\section*{Quick Reference Card}
\addcontentsline{toc}{section}{Quick Reference Card}
% ============================================================

\subsection*{TwinCAT Function Series}

\begin{table}[h]
\centering
\begin{tabular}{@{}ll@{}}
\toprule
\textbf{Series} & \textbf{Category} \\
\midrule
TF1xxx & System \\
TF2xxx & HMI \\
TF3xxx & Measurement \\
TF4xxx & Controller \\
TF5xxx & Motion Control \\
TF6xxx & Connectivity \\
TF7xxx & Vision \\
TF8xxx & Industry Specific \\
\bottomrule
\end{tabular}
\end{table}

\subsection*{Modbus Function Blocks}

\begin{table}[h]
\centering
\begin{tabular}{@{}ll@{}}
\toprule
\textbf{Function Block} & \textbf{Purpose} \\
\midrule
FB\_MBReadCoils & Read digital outputs (coils) \\
FB\_MBReadInputs & Read digital inputs \\
FB\_MBReadRegs & Read holding registers \\
FB\_MBReadInputRegs & Read input registers \\
FB\_MBWriteCoils & Write multiple coils \\
FB\_MBWriteSingleRegister & Write single register \\
FB\_MBWriteRegs & Write multiple registers \\
\bottomrule
\end{tabular}
\end{table}

\subsection*{Asynchronous Pattern}

\begin{lstlisting}
// Standard pattern for all async Modbus operations
CASE eState OF
    WAIT:
        IF bTrigger THEN eState := TRIGGER; END_IF
        
    TRIGGER:
        fbModbus(bExecute := TRUE, ...);
        eState := BUSY;
        
    BUSY:
        fbModbus(bExecute := FALSE);
        IF NOT fbModbus.bBusy THEN
            // Check bError, process results
            eState := WAIT;
        END_IF
END_CASE
\end{lstlisting}

\subsection*{Data Type Conversions}

\begin{lstlisting}
// INT to WORD (for writing)
wValue := INT_TO_WORD(nIntValue);

// WORD to INT (for reading)
nIntValue := WORD_TO_INT(wValue);

// BOOL to BYTE (for coils)
byValue := BOOL_TO_BYTE(bValue);
\end{lstlisting}

\vfill

\begin{center}
\textit{Last Updated: \today}
\end{center}

\end{document}
